\pendiente{Este capítulo se redacta con los resultados de la evaluación (semanas del 19-oct al 1-nov). La estructura prevista es la siguiente.}

\section{Corpus y conjunto de consultas}
\pendiente{Estadísticas del corpus congelado: ofertas por fuente, fragmentos por tipo, tasa de análisis estructurado fallido. Descripción del conjunto de consultas por categoría e idioma.}

\section{Validación del juez de relevancia}
\pendiente{Tabla: kappa binario y ponderado (con IC 95\,\%) del juez local y del juez final frente a los juicios humanos; matriz de confusión.}

\section{Calidad de recuperación (PE1)}
\pendiente{Tabla: P@5, P@10, nDCG@10, MRR y Recall@10 por sistema (ILIKE, BM25, pgvector exacto, pgvector HNSW, Qdrant HNSW), con significancia frente a cada línea base. Figura: nDCG@10 por categoría de consulta. Tabla: resultados del subconjunto en español. Exactitud del enrutamiento en consultas de agregación.}

\section{Fidelidad de la generación (PE2)}
\pendiente{Tabla: fidelidad y relevancia de la respuesta por categoría; análisis de casos de baja fidelidad.}

\section{Latencia del pipeline (PE3)}
\pendiente{Tabla: p50 y p95 por etapa (\textit{embedding}, búsqueda, construcción del contexto, TTFT, total) con cada almacén vectorial.}

\section{pgvector frente a Qdrant (PE4)}

\subsection{Resultado preliminar sobre el corpus real}

La Tabla~\ref{tab:res-ann-real} resume una primera comparación sobre el corpus de evaluación (4.013 fragmentos). Se usaron 200 consultas formadas por títulos de ofertas del corpus y 3 repeticiones; la referencia es la búsqueda exacta en pgvector. La latencia se mide desde el cliente y corresponde a la búsqueda de 10 vecinos, sin el \textit{embedding} de la consulta.

\begin{table}[htbp!]
    \centering
    \caption{Recall@10 y latencia de búsqueda sobre el corpus real (resultado preliminar)}
    \label{tab:res-ann-real}
    \begin{tabular}{|l|c|c|c|}
        \hline
        \textbf{Configuración} & \textbf{Recall@10} & \textbf{p50 (ms)} & \textbf{p95 (ms)} \\
        \hline
        pgvector, exacta & 1{,}000 & 24{,}3 & 27{,}2 \\
        \hline
        Qdrant, exacta & 0{,}999 & 3{,}9 & 4{,}8 \\
        \hline
        pgvector, plan automático & 1{,}000 & 23{,}6 & 25{,}3 \\
        \hline
        pgvector, HNSW $ef=16$ & 0{,}976 & 2{,}0 & 2{,}7 \\
        \hline
        Qdrant, HNSW $ef=16$ & 0{,}925 & 3{,}0 & 3{,}7 \\
        \hline
        pgvector, HNSW $ef=64$ & 0{,}995 & 2{,}2 & 2{,}8 \\
        \hline
        Qdrant, HNSW $ef=64$ & 0{,}988 & 3{,}1 & 3{,}6 \\
        \hline
        pgvector, HNSW $ef=256$ & 0{,}999 & 3{,}2 & 3{,}8 \\
        \hline
        Qdrant, HNSW $ef=256$ & 0{,}998 & 3{,}5 & 4{,}1 \\
        \hline
    \end{tabular}
\end{table}

A esta escala se observan tres hechos que el benchmark completo deberá contrastar:
\begin{itemize}
    \item Con el plan automático, pgvector recorre la tabla completa (Sección~\ref{sec:art-vectorial}), por lo que la aplicación servida desde pgvector tarda unas ocho veces más que servida desde Qdrant.
    \item Forzando el índice, pgvector es ligeramente más rápido que Qdrant y alcanza un \textit{recall} mayor con el mismo $ef_{search}$. A este tamaño, la latencia de Qdrant está dominada por la sobrecarga del cliente y del protocolo.
    \item Las diferencias de \textit{recall} entre los dos motores en modo exacto se deben a empates entre fragmentos idénticos.
\end{itemize}

\pendiente{Figuras: curvas \textit{recall}@10 frente a latencia por escala (corpus real, 100k, 500k) y por selectividad del filtro (sin filtro, 50\,\%, 10\,\%, 1\,\%). Tabla: tiempo de construcción, tamaño del índice, memoria pico y QPS.}

\section{Discusión}
\pendiente{Respuesta a cada pregunta de evaluación, relación con la literatura revisada en el Capítulo~\ref{cap:estado-arte} y limitaciones observadas.}
